

\newcommand{\mat}[1]{\begin{bmatrix}#1\end{bmatrix}}
\newcommand{\vet}[1]{\begin{bmatrix}#1\end{bmatrix}}

\begin{question}
  \begin{block}{}
    Considerando a equação
    \[
      x^2+4xy+y^2-2x-2y-3=0
    \]
  \end{block}

  \begin{enumerate}
    \setlength{\itemsep}{0.45cm}
    \item escreva a equação na forma matricial
    \item diagonalize matriz da parte quadrática
    \item faça a mudança de variáveis

      (apresente a transformação e transforme a equação)
    \item identifique a cônica
  \end{enumerate}
\end{question}

\begin{resolution}{Equação Matricial}
  \textbf{Equação}
  \[
    x^2+4xy+y^2-2x-2y-3=0
  \]
  \[
    ax^2+2bxy+cy^2+dx+ey+f=0
  \]

  \textbf{Equação na forma matricial}
  \[
    v^tAv+B^tv+f=0
  \]

  \[
    v=\mat{x\\y}
    \qquad
    A=\mat{a&b\\b&c}
     =\mat{1&2\\2&1}
    \qquad
    B=\mat{d\\e}
     =\mat{-2\\-2}
    \qquad
    f=-3
  \]
\end{resolution}

\begin{resolution}{Diagonalizando A -- Polinômio Característico}
  \textbf{Polinômio característico}
  \[
    p(\lambda)=\det(A-\lambda I)
  \]
  \[
    =\det\left(
      \mat{1&2\\2&1}
      -\lambda\mat{1&0\\0&1}
    \right)
  \]
  \[
    =\begin{vmatrix}
      1-\lambda&2\\
      2&1-\lambda
    \end{vmatrix}
  \]
  \[
    =(1-\lambda)^2-4
  \]
  \[
    =1-2\lambda+\lambda^2-4
  \]
  \[
    =\lambda^2-2\lambda-3
  \]
\end{resolution}

\begin{resolution}{Diagonalizando A -- Equação Característica}
  \textbf{Equação característica}
  \[
    p(\lambda)=\lambda^2-2\lambda-3=0
  \]
  \[
    \Delta=b^2-4ac=(-2)^2-4\times1(-3)=16
    \qquad
    \sqrt{\Delta}=4
  \]
  \[
    \lambda_1=\frac{-b+\sqrt{\Delta}}{2a}
    =\frac{2+4}{2}=3
  \]
  \[
    \lambda_2=\frac{-b-\sqrt{\Delta}}{2a}
    =\frac{2-4}{2}=-1
  \]

  \textbf{Autovalores}
  \[
    \lambda_1=3
    \qquad
    \lambda_2=-1
  \]
\end{resolution}

\begin{resolution}{Autovetores Associados ao Autovalor $\lambda_1=3$}
  \begin{columns}[T]
    \begin{column}{0.58\textwidth}
      \[
        (A-3I)v_1=\vec{0}
      \]
      \[
        \left(
        \mat{1&2\\2&1}
        -3\mat{1&0\\0&1}
        \right)
        \mat{x\\y}
        =
        \mat{0\\0}
      \]
      \[
        \mat{-2&2\\2&-2}
        \mat{x\\y}
        =
        \mat{0\\0}
      \]
    \end{column}
    \begin{column}{0.37\textwidth}
      \[
        -x+y=0
      \]
      \[
        y=x
      \]
      Escolhendo $x=1$
      \[
        v_1=\mat{1\\1}
      \]
    \end{column}
  \end{columns}
\end{resolution}

\begin{resolution}{Autovetores Associados ao Autovalor $\lambda_2=-1$}
  \begin{columns}[T]
    \begin{column}{0.58\textwidth}
      \[
        (A-(-1)I)v_2=\vec{0}
      \]
      \[
        \left(
        \mat{1&2\\2&1}
        +\mat{1&0\\0&1}
        \right)
        \mat{x\\y}
        =
        \mat{0\\0}
      \]
      \[
        \mat{2&2\\2&2}
        \mat{x\\y}
        =
        \mat{0\\0}
      \]
    \end{column}
    \begin{column}{0.37\textwidth}
      \[
        2x+2y=0
      \]
      \[
        y=-x
      \]
      Escolhendo $x=1$
      \[
        v_2=\mat{1\\-1}
      \]
    \end{column}
  \end{columns}
\end{resolution}

\begin{resolution}{Matriz Ortogonal}
  \textbf{Normalizando o autovetor $v_1$}
  \[
    \|v_1\|=\sqrt{1^2+1^2}=\sqrt{2}
    \qquad
    u_1=\frac{v_1}{\|v_1\|}
    =\frac{1}{\sqrt{2}}\mat{1\\1}
  \]

  \textbf{Normalizando o autovetor $v_2$}
  \[
    \|v_2\|=\sqrt{1^2+(-1)^2}=\sqrt{2}
    \qquad
    u_2=\frac{v_2}{\|v_2\|}
    =\frac{1}{\sqrt{2}}\mat{1\\-1}
  \]

  \textbf{Matriz ortogonal $P$ e matriz diagonal $D$}
  \[
    P=[u_1\ u_2]
    =\frac{1}{\sqrt{2}}
      \mat{1&1\\1&-1}
    \qquad
    D=\operatorname{diag}(\lambda_1,\lambda_2)
    =\mat{3&0\\0&-1}
  \]
\end{resolution}

\begin{resolution}{Mudança de Variáveis}
  \textbf{Mudança de variáveis $v=Pw$}
  \[
    \mat{x\\y}
    =
    \frac{1}{\sqrt{2}}
    \mat{1&1\\1&-1}
    \mat{\mu\\\nu}
    =
    \frac{1}{\sqrt{2}}
    \mat{\mu+\nu\\\mu-\nu}
    =
    \mat{\dfrac{\mu+\nu}{\sqrt{2}}\\[2mm]
         \dfrac{\mu-\nu}{\sqrt{2}}}
  \]
  \[
    x=\frac{\mu+\nu}{\sqrt{2}}
    \qquad
    y=\frac{\mu-\nu}{\sqrt{2}}
  \]
\end{resolution}

\begin{resolution}{Transformando as Matrizes}
  \[
    D=P^tAP=\operatorname{diag}(\lambda_1,\lambda_2)
    =\mat{3&0\\0&-1}
  \]

  \[
    E=P^tB
    =\frac{1}{\sqrt{2}}
      \mat{1&1\\1&-1}
      \mat{-2\\-2}
    =\frac{1}{\sqrt{2}}
      \mat{-4\\0}
    =\mat{-2\sqrt{2}\\0}
  \]
\end{resolution}

\begin{resolution}{Equação Transformada}
  \[
    w^tDw+E^tw+f=0
  \]

  \[
    \begin{bmatrix}\mu&\nu\end{bmatrix}
    \mat{3&0\\0&-1}
    \mat{\mu\\\nu}
    +
    \begin{bmatrix}-2\sqrt{2}&0\end{bmatrix}
    \mat{\mu\\\nu}
    -3=0
  \]

  \[
    3\mu^2-\nu^2-2\sqrt{2}\mu-3=0
  \]
\end{resolution}

\begin{resolution}{Completando Quadrados}
  \[
    3\mu^2-\nu^2-2\sqrt{2}\mu-3=0
  \]
  \[
    (3\mu^2-2\sqrt{2}\mu)-\nu^2-3=0
  \]
  \[
    3\left(\mu^2-\frac{2\sqrt{2}}{3}\mu\right)
    -\nu^2-3=0
  \]

  \[
    \begin{aligned}
      \mu^2-\frac{2\sqrt{2}}{3}\mu
      &=
      \mu^2-\frac{2\sqrt{2}}{3}\mu
      +\left(\frac{\sqrt{2}}{3}\right)^2
      -\left(\frac{\sqrt{2}}{3}\right)^2\\
      &=
      \left(\mu-\frac{\sqrt{2}}{3}\right)^2-\frac{2}{9}
    \end{aligned}
  \]
\end{resolution}

\begin{resolution}{Completando Quadrados}
  \begin{columns}[T]
    \begin{column}{0.49\textwidth}
      \[
        3\left(\mu^2-\frac{2\sqrt{2}}{3}\mu\right)
        -\nu^2-3=0
      \]
      \[
        3\left[
        \left(\mu-\frac{\sqrt{2}}{3}\right)^2
        -\frac{2}{9}
        \right]-\nu^2=3
      \]
      \[
        3\left(\mu-\frac{\sqrt{2}}{3}\right)^2
        -\frac{2}{3}-\nu^2=3
      \]
      \[
        3\left(\mu-\frac{\sqrt{2}}{3}\right)^2
        -\nu^2=\frac{11}{3}
      \]
    \end{column}
    \begin{column}{0.49\textwidth}
      \[
        \frac{
        3\left(\mu-\frac{\sqrt{2}}{3}\right)^2
        }{11/3}
        -
        \frac{\nu^2}{11/3}=1
      \]
      \[
        \frac{
        \left(\mu-\frac{\sqrt{2}}{3}\right)^2
        }{11/9}
        -
        \frac{\nu^2}{11/3}=1
      \]
    \end{column}
  \end{columns}
\end{resolution}

\begin{resolution}{Identificação}
  \textbf{A cônica}
  \[
    \frac{\left(\mu-\frac{\sqrt{2}}{3}\right)^2}{11/9}
    -\frac{\nu^2}{11/3}=1
  \]

  é uma hipérbole com

  \textbf{Centro}
  \[
    C=\left(\frac{\sqrt{2}}{3},0\right)
    \approx(0{,}816;0)
  \]

  \textbf{Parâmetros}
  \[
    a=\sqrt{\frac{11}{9}}=\frac{\sqrt{11}}{3}
    \approx1{,}105
    \qquad
    b=\sqrt{\frac{11}{3}}
    \approx1{,}915
  \]
\end{resolution}


%------------------------------------------------------------------------------%
\begin{resolution}{Esboço}
  \centering
  \includegraphics{I-08-Rotacao-ex-2}
\end{resolution}
%------------------------------------------------------------------------------%
